\section{Introduction and motivation}
Our question is whether Linux queue management can protect a small
request-response flow when bulk TCP traffic competes for the same bottleneck.
We compare FIFO, FQ-CoDel, CAKE best-effort, and CAKE with bounded classification.
The goal is to reduce added foreground p99 latency by at least 50\% relative
to FIFO while retaining at least 90\% of FIFO's background goodput.
These thresholds are our feasibility criteria, not established results.

We proposed a three-node Linux testbed, repeated randomized trials, and a
comparison of latency against throughput. Hardware availability changed the
endpoint plan: we now intend to use two MacBooks with a Linux NucBox router.
We also developed a namespace path for controlled software checks and a local
simulation for exercising the analysis. These paths answer different questions.
The physical experiment will test the available equipment; the simulation
cannot tell us which Linux queue policy performs best.

Our preliminary namespace pilot gives a reason to evaluate both goals together.
FQ-CoDel and both CAKE configurations had much lower observed added p99 than
FIFO. Both CAKE variants also retained less than 90\% of FIFO goodput. With only
two paired blocks, generator omissions and unresolved timing checks, we treat
this as a diagnostic result that guides the next experiment.

On November 26 we also ran the workload on this Mac through loopback and
integrated Chris's RPC implementation into the trial runner. The foreground
phase retained 20,000 scheduled slots, including 53 that were not sent. All
20,000 calls in the separate RPC phases completed. This checks the local
collection path while leaving the physical queue-policy question open.
